\section{Introduction}\label{intro:section}

The Kodaira dimension measures the growth of pluricanonical forms.
For a surjective morphism $f:X\to Y$ of smooth projective complex varieties
with connected fibers, Iitaka's subadditivity problem asks whether
\[
                  \kappa(X)\geq\kappa(F)+\kappa(Y),
\]
where $F$ is a general fiber. The logarithmic version permits poles
along reduced boundary divisors and therefore applies to open
varieties. We prove a common generalization in which rational boundary
coefficients and multiple fibers both contribute to the inequality.

Iitaka's divisor-dimension and logarithmic theories provide the
language for this question \cite{IitakaOriginal,IitakaLog}; an
explicit logarithmic subadditivity conjecture appears in
\cite[\S3]{IitakaClassification}.
The lower bound requires global pluricanonical forms that retain
the growth of both the fibers and the base. Direct-image positivity
established central cases: Kawamata treated curve bases and fibers
admitting good minimal models, Viehweg developed weak positivity
and addition over a base of general type, and Koll\'ar proved
addition for general-type fibers
\cite{KawamataCurves,KawamataKodaira,ViehwegAdditivity,Kollar87}.
Weak positivity controls symmetric powers after small positive
ample twists; a general-type base provides enough canonical
positivity to absorb them.

Other advances exploit dimension or the geometry of the base:
Birkar treated total dimension at most six, Cao--P\u{a}un and
Hacon--Popa--Schnell treated abelian or maximal-Albanese-dimension
bases, and Cao treated projective klt pairs over surfaces
\cite{BirkarIitakaSix,CaoPaunAbelian,HaconPopaSchnellAbelian,CaoSurfaces}.
Thus the established cases draw positivity from different parts of
the fibration, rather than imposing one uniform condition on it.

For logarithmic pairs, Kov\'acs--Patakfalvi proved the projective
log-general-type-fiber case \cite[Theorem~9.6]{KovacsPatakfalvi},
while Hashizume proved reduced-SNC subadditivity for fibers with
abundant log canonical divisor \cite[Theorem~1.2]{Hashizume}.
Maehara had established logarithmic addition over a base of log
general type; Fujino gives a proof through twisted weak positivity
\cite[Theorem~1.9]{FujinoWeakPositivity}.
Campana proved addition over an orbifold base of general type, and
Wang proved klt K\"ahler subadditivity over complex tori, including
an inf-multiplicity refinement
\cite[Theorem~6.3]{CampanaOrbifolds}
\cite[Main Theorem~A and Theorem~E]{WangKahlerIitaka}.
The present argument uses general-type addition only after a relative
Iitaka reduction; the original fiber need not be of general type
or admit a good minimal model.

Earlier preprints also state general ordinary inequalities.
Tsuji states ordinary subadditivity following a direct-image generation
theorem \cite[Theorems~1.9 and~1.13]{TsujiDirectImages}; his
Remark~1.14 distinguishes an unproved variation refinement.
Maehara states an ordinary determinant inequality involving variation
\cite[\S7, Theorem~8]{MaeharaIitaka}. This later preprint claim is
distinct from his logarithmic general-type-base theorem above.
Neither preprint statement is a proof input here.

\subsection{The orbifold base and the main theorem}

Multiple fibers carry information that the ordinary base omits.
Campana's orbifold base records it through inf multiplicities, and
his birational invariant gives the corresponding form of Iitaka's
conjecture \cite{Campana2004,CampanaOrbifolds}. We specify that
invariant before stating the theorem.

\label{setup:orbifold-base}
A fibration is a proper surjective holomorphic map with connected
fibers. A rational SNC boundary on a smooth manifold is a finite
sum $\Delta=\sum_i\delta_i\Delta_i$, with
$\delta_i\in\mathbf Q\cap[0,1]$ and simple normal crossing support.
Write $\Delta_E$ for the coefficient of a source prime $E$. Set
\[
 m_\Delta(E)=\frac1{1-\Delta_E},
 \qquad \frac1{0}=\infty,
\]
where $\Delta_E=0$ off the boundary. If the base of
$f:(X,\Delta)\to Y$ is smooth, define
\begin{equation}\label{setup:inf-boundary}
\begin{gathered}
 m(f,\Delta;D)=\min_{E\mapsto D}
     \bigl\{\operatorname{ord}_E(f^*D)m_\Delta(E)\bigr\},
 \qquad
 B(f,\Delta)=\sum_D\left(1-\frac1{m(f,\Delta;D)}\right)D.
\end{gathered}
\end{equation}
Here $D$ ranges over base prime divisors, $E$ over source primes
dominating $D$, and $1/\infty=0$. The minimum is nonempty and
only finitely many coefficients of $B(f,\Delta)$ are nonzero.
There is no divisibility condition on these multiplicities.

For smooth sources and bases, an elementary change of model is a
commutative diagram of fibrations
\[
\begin{tikzcd}
 (X',\Delta') \arrow[r,"u"] \arrow[d,"f'"'] &
 (X,\Delta) \arrow[d,"f"]\\
 Y' \arrow[r,"v"'] & Y
\end{tikzcd}
\]
where $u,v$ are proper bimeromorphic holomorphic maps,
$u_*\Delta'=\Delta$, and both boundaries are rational SNC
boundaries in $[0,1]$. We require, for every prime $D$ on $X$
and prime $E$ on $X'$ with $t=\operatorname{ord}_E(u^*D)>0$,
\begin{equation}\label{setup:orbifold-morphism}
                     t\,m_{\Delta'}(E)\geq m_\Delta(D).
\end{equation}
We permit $\infty\geq\infty$ and impose no condition for $t=0$.
Write $f'\sim f$ for the equivalence relation generated by these
diagrams, allowing arrows in either direction. Define
\begin{equation}\label{setup:invariant-base}
 \kappa(f\mid\Delta)=
 \inf_{f':(X',\Delta')\to Y'\sim f}
       \kappa\bigl(Y',K_{Y'}+B(f',\Delta')\bigr).
\end{equation}
For a normal singular base $Y$, resolve $Y$, resolve the main
component of the fiber product with $X$, and give the resulting
smooth source the strict transform of $\Delta$ plus the reduced
exceptional divisor over $X$. After resolving this boundary, apply
the preceding definition. It is independent of the chosen resolution
\cite[Definitions~2.1, 2.3, 3.2, 4.7, and~4.10;
Remark~4.13 and Corollary~4.14]{CampanaOrbifolds}.

For a rational line bundle $L$, $\kappa(Z,L)$ is the maximal image
dimension of its complete linear systems in sufficiently divisible
positive degrees, and is $-\infty$ if all these systems are empty.
A point has Kodaira dimension zero, and
$(-\infty)+a=-\infty$, including $a=-\infty$.
Fujiki class $\mathcal C$ consists of compact complex spaces
bimeromorphic to compact K\"ahler spaces.

\begin{theorem}[Orbifold subadditivity]\label{main:orbifold}
Let $X$ be a smooth compact connected complex manifold in Fujiki
class $\mathcal C$, let $\Delta$ be a rational SNC boundary, and
let $f:X\to Y$ be a surjective holomorphic map with connected
fibers onto a normal compact irreducible complex space.
For a very general smooth fiber $F$, put $\Delta_F=\Delta|_F$.
Then
\begin{equation}\label{main:inequality}
 \kappa(X,K_X+\Delta)
 \geq\kappa(F,K_F+\Delta_F)+\kappa(f\mid\Delta).
\end{equation}
\end{theorem}

This is a positive resolution of Campana's orbifold Iitaka
conjecture in its rational SNC, Fujiki-class-$\mathcal C$
formulation, including coefficient one
\cite[Conjecture~6.1 and Remark~6.2]{CampanaOrbifolds}.
Here a very general fiber is chosen outside a countable union of
proper closed analytic subsets, including the critical values and
the exceptional images of boundary strata.

For reduced SNC divisors $D_X,D_Y$ satisfying
$\operatorname{Supp}(f^*D_Y)\subseteq\operatorname{Supp}(D_X)$,
the invariant base term is at least $\kappa(Y,K_Y+D_Y)$.
Theorem~\ref{main:orbifold} therefore gives
\[
 \kappa(X,K_X+D_X)
 \geq\kappa(F,K_F+D_X|_F)+\kappa(Y,K_Y+D_Y)
\]
for fibrations between smooth compact class-$\mathcal C$ manifolds.
For a smooth quasi-projective variety $U$, its logarithmic Kodaira
dimension is $\bar\kappa(U)=\kappa(X,K_X+D_X)$ on a smooth
projective compactification with reduced SNC boundary
$D_X=X\setminus U$. Compatible compactifications give logarithmic
subadditivity for dominant morphisms of smooth quasi-projective
complex varieties with geometrically connected general fiber.
Empty boundaries give ordinary subadditivity. The projective
ordinary and logarithmic statements also hold over algebraically
closed fields of characteristic zero, using geometric generic fibers.

\subsection{The reduction and the two positivity inputs}

The proof first replaces the original fiber by a general-type
object without losing its complete pluricanonical systems.
Assume the two terms on the right of Theorem~\ref{main:orbifold}
are nonnegative, and put $k=\kappa(F,K_F+\Delta_F)$.
A relative Iitaka map gives, after the model changes of Section~\ref{setup:section},
\[
 X\xrightarrow{g}W\xrightarrow{h}Y,
 \qquad \dim W_y=k,
 \qquad\kappa(G,K_G+\Delta_G)=0
\]
for a very general $g$-fiber $G$.

The canonical-bundle-formula strategy separates the contribution
of singular fibers from a moduli contribution on $W$.
Fujino--Mori developed a higher-dimensional logarithmic formula,
and Ambro organized its discriminant and moduli parts birationally
\cite{FujinoMori,AmbroBoundary,AmbroModuli}.
Our use of this strategy requires an exact comparison on the fixed
original source. Section~\ref{setup:section} obtains it from the orders of relative
pluriforms. A sufficiently divisible pluricanonical system on $G$
is one-dimensional. Taking a root of its generator on a cyclic
cover gives a differential form, and the line it spans belongs to
the cohomology of that cover. Tensor descent and extension across
the boundary produce a rational line bundle $M$ on $W$.
The orders of the same form determine a boundary $T$ such that
\[
 B(g,\Delta)\leq T\leq1,
 \qquad
 H^0\bigl(X,m(K_X+\Delta)\bigr)
 \simeq H^0\bigl(W,m(K_W+T+M)\bigr)
\]
in sufficiently divisible degrees. The comparison also holds
relatively over $Y$. Thus $K_W+T+M$ is big on the
$k$-dimensional fibers of $h$.

Two inputs turn this fiberwise bigness into the required global
lower bound. The first concerns the line $M$: the period map,
which records how the Hodge structures of the covers vary, supplies
a second fibration $p:W\to S$, with $S$ smooth projective, and
a nef rational line bundle $P$ such that
\[
                      M=p^*P,\qquad K_S+aP\ \text{is big}
\]
for some positive rational $a$. The maps $h$ and $p$ have different
jobs, as Figure~\ref{fig:iitaka-period} shows. The second input is
log-general-type addition: a log canonical divisor that is big on
the fibers of $h$ has Kodaira dimension at least their dimension
plus the orbifold contribution of $Y$.

\begin{figure}[ht]
\centering
\begin{tikzcd}[column sep=large,row sep=large]
 X \arrow[r,"g"] & W \arrow[r,"h"] \arrow[d,"p"'] & Y\\
 & S &
\end{tikzcd}
\caption{The relative Iitaka reduction factors the original fibration
through $W$. The map $h$ retains the original base and has
$k$-dimensional general fibers. The independent map $p$ expresses
$M$ as the pullback of a nef line $P$ on a projective base where
$K_S+aP$ is big. No map between $Y$ and $S$ is asserted.}
\label{fig:iitaka-period}
\end{figure}


Section~\ref{addition:section} states both inputs precisely and proves the intervening
addition principle. Write $D_W=K_W+T$. For a fixed rational $c<1$
close to one, general-type addition gives sufficiently many ratios
of sections of $D_W+cM+\eta p^*A$, where $A$ is ample on $S$
and $\eta>0$ is rational. Weak positivity then supplies one fixed
section of $D_W+aM-\lambda p^*A$ for suitable $a>1$ and
$\lambda>0$. A positive rational combination cancels the ample
terms and gives coefficient one on $M$. Multiplication by powers
of that fixed section preserves the ratios of the first system.
This is an exact section argument; it requires neither abundance
of $P$ nor a limiting assertion for Kodaira dimensions.

The full proofs of the inputs follow this calculation.
Section~\ref{period:section} proves adjoint positivity of the Hodge line. It combines
the curvature and boundary theory of Griffiths, Schmid, and
Cattani--Kaplan--Schmid with a projective period quotient
\cite{GriffithsPeriodsIII,Schmid,CKS}.
The construction retains suitable rational adjoint factors of the
ambient integral variation, associated with its monodromy, from which
a positive power of the chosen line can be recovered. Their full
periods give the quotient, using Villadsen's class-$\mathcal C$
compactification theorem and the period-quotient approach of
Bakker--Brunebarbe--Tsimerman and Sommese; curvature and a volume
argument establish its projectivity
\cite{BakkerBrunebarbeTsimerman,SommesePeriodMapping,VilladsenHodgeKahler}.

On this quotient $S$, the retained full period map is generically
immersive, but the map remembering only the chosen line may have
smaller rank. Curvature identifies the latter rank with the numerical
dimension $\nu(P)$, measured by its nonzero intersection powers.
Let $D$ be the reduced boundary where the retained variation has
nonidentity local monodromy. Brunebarbe--Cadorel gives $K_S+D$ big
\cite{BrunebarbeCadorel}. Using Andr\'e's monodromy normality and
Bakker--Tsimerman's Ax--Schanuel theorem to analyze fibers of the
line map, the boundary argument proves
$\nu(P|_{D_i})<\nu(P)$ for every component $D_i$ of $D$, including
finite nonidentity monodromy
\cite{AndreMonodromy,BakkerTsimermanAxSchanuel}.
The boundary loss in the asymptotic section count consequently has
lower degree in the coefficient of $P$ than the available sections
on $S$. Lemma~\ref{period:subtract} turns this comparison into
bigness of $K_S+aP$.

Section~\ref{generaltype:section} proves log-general-type addition in class $\mathcal C$.
Kov\'acs--Patakfalvi's stable-family positivity applies to projective
families. We construct such an auxiliary family and descend its
relative sections to the original class-$\mathcal C$ base.
Fujino's weak positivity supplies the other direct-image input
used in Section~\ref{addition:section} \cite{KovacsPatakfalvi,FujinoWeakPositivity}.
Section~\ref{application:section} then returns once to the original fibration, applies the
addition principle to $(W,T,M)$, and derives the ordinary,
logarithmic, characteristic-zero, and core consequences. The core
and the relevant specialness of Kodaira-zero fibers are Campana's
constructions and results \cite{Campana2004,CampanaOrbifolds}.

\subsection{The reduced-boundary refinement}

For a rational boundary, the root form above spans one character
subspace of the cyclic cover's highest Hodge space; the whole
highest space can be larger. For projective reduced-boundary pairs
of logarithmic Kodaira dimension zero, Section~\ref{cyc:section} proves more:
the minimal root cover has a one-dimensional entire logarithmic
top-form space. Mixed-Hodge extension and semipositivity in the
forms of Fujino--Fujisawa control the coefficient-one boundary
\cite{FujinoFujisawa,FujinoFujisawaVMHS}.
We identify the resulting line with the moduli divisor in the
projective subpair setting of Bakker--Filipazzi--Mauri--Tsimerman
\cite[Theorem~6.28]{BFMT}. This identification does not use their
separate semiampleness theorem, whose hypotheses include
effectiveness over the generic point \cite[Theorem~1.5]{BFMT}.
Section~\ref{cyc:section} also gives the exact divisorial normalization.
Appendix~\ref{p1logcbfalt:section} collects complementary projective
calculations by valuations and toroidal models, together with the
comparison on higher models.

The logarithmic lower bound, adjoint positivity, and this
reduced-boundary refinement supply the stated inputs to the
whole-fiber-variation companion \cite{IitakaVariationCompanion}. The reverse-inequality companion
\cite{IitakaAdditivityCompanion} uses the lower bound only for additivity; its upper theorem has
additional smoothness hypotheses on all boundary strata.
Neither companion is used in the proof of Theorem~\ref{main:orbifold}.
The separate article on projective Hodge lines
\cite{IitakaProjectiveCompanion} gives a complete
ordinary proof through rationally polarized integral variations,
together with the distinct integral and analytic boundary arguments.
